\documentclass[a4paper,12pt]{extarticle}
\usepackage{geometry}
\usepackage{fontspec}
\setmainfont{Times New Roman}
\usepackage[main=russian,english]{babel}
\usepackage{csquotes}
\usepackage{amsmath}
\usepackage{amsthm}
\usepackage{amssymb}
\usepackage{fancyhdr}
\usepackage{setspace}
\usepackage{graphicx}
\usepackage{colortbl}
\usepackage{tikz}
\usepackage{pgf}
\usepackage{subcaption}
\usepackage{listings}
\usepackage{indentfirst}
\usepackage[
backend=biber,
style=numeric,
maxbibnames=99
]{biblatex}
\addbibresource{refs.bib}
\usepackage[colorlinks,citecolor=blue,linkcolor=blue,bookmarks=false,hypertexnames=true, urlcolor=blue]{hyperref} 
\usepackage{indentfirst}
\usepackage{mathtools}
\usepackage{booktabs}
\usepackage[flushleft]{threeparttable}
\usepackage{tablefootnote}

\usepackage{chngcntr} % нумерация графиков и таблиц по секциям
\counterwithin{table}{section}
\counterwithin{figure}{section}

\graphicspath{{graphics/}}%путь к рисункам

\makeatletter
% \renewcommand{\@biblabel}[1]{#1.} % Заменяем библиографию с квадратных скобок на точку:
\makeatother

\geometry{left=2.5cm}% левое поле
\geometry{right=1.0cm}% правое поле
\geometry{top=2.0cm}% верхнее поле
\geometry{bottom=2.0cm}% нижнее поле
\setlength{\parindent}{1.25cm}
\renewcommand{\baselinestretch}{1.5} % междустрочный интервал


\newcommand{\bibref}[3]{\hyperlink{#1}{#2 (#3)}} % biblabel, authors, year
\addto\captionsrussian{\def\refname{Список литературы (или источников)}} 

\renewcommand{\theenumi}{\arabic{enumi}}% Меняем везде перечисления на цифра.цифра
\renewcommand{\labelenumi}{\arabic{enumi}}% Меняем везде перечисления на цифра.цифра
\renewcommand{\theenumii}{.\arabic{enumii}}% Меняем везде перечисления на цифра.цифра
\renewcommand{\labelenumii}{\arabic{enumi}.\arabic{enumii}.}% Меняем везде перечисления на цифра.цифра
\renewcommand{\theenumiii}{.\arabic{enumiii}}% Меняем везде перечисления на цифра.цифра
\renewcommand{\labelenumiii}{\arabic{enumi}.\arabic{enumii}.\arabic{enumiii}.}% Меняем везде перечисления на цифра.цифра

\begin{document}
\input{title_kr}% это титульный лист - выберите подходящий вам из имеющихся в проекте вариантов (kr - курсовая работа у 3 курса, vkr - выпускная квалификационная работа у 4 курса)
\newpage
\setcounter{page}{2}

{
	\hypersetup{linkcolor=black}
	\tableofcontents
}

\newpage

\newpage
\section*{Аннотация}   % this is how to use russian
При проектировании высокопроизводительных и распределенных систем одной из важнейших задач является выбор и настройка системы выдачи заданий (балансировщик нагрузки). При использовании гомогенных вычислительных узлов в кластере можно использовать готовые решения (slurm, torque и др.), но в гетерогенных системах готовые решения не будут работать из коробки. В распределенных системах из-за большой гетерогенности и других ограничений может возникнуть потребность разработки оригинального планировщика заданий. В распределенной вычислительной системе в качестве части ресурсов могут быть использованы облачные ресурсы или ресурсы кластеров. Уже были разработаны решения (к примеру, CluBoRun в грид-системе на платформе BOINC) по использованию ресурсов кластера в фоновом режиме. Предлагается разработать /модифицировать систему выдачи заданий для грид-системы на платформе BOINC, в которой присутствуют сегменты кластеров.

\addcontentsline{toc}{section}{Аннотация}

\section*{Ключевые слова}
Распределенные системы, балансировка нагрузки
\pagebreak

\section{Введение} 
\subsection{Описание предметной области}
Предметная область проекта охватывает распределенные системы с гетерогенной инфраструктурой с основными компонентами:
\begin{itemize}
    \item BOINC (Berkeley Open Infrastructure for Network Computing) — платформа для добровольных вычислений, объединяющая ресурсы разнородных устройств (серверов, кластеров, персональных компьютеров) в единую сеть
    \item Кластеры — узлы, интегрированные в распределенную систему с целью повышения доступности
    \item Гетерогенные ресурсы — устройства с различными характеристиками: вычислительной мощностью, объемом памяти, доступностью и другими
    \item Балансировщик нагрузки — компонент, отвечающий за распределение вычислительных задач между узлами с учетом их возможностей и текущих состояний
\end{itemize}
Также предметная область имеет следующие особенности:
\begin{itemize}
    \item Отсутствие централизованного контроля за узлами системы
    \item Разнородность ресурсов каждого из узлов
    \item Различные режимы работы узлов (не все устройства доступны на протяжении 100\% времени)
    \item Различный объем задач, время выполнения которых может варьироваться от нескольких минут до нескольких дней
    \item Необходимость минимизации накладных расходов для балансировщика
\end{itemize}

\subsection{Релевантность и значимость}
Актуальность данного проекта обусловлена множеством существующих на данный момент проблем, таких как:
\begin{enumerate}
    \item Современные вычислительные задачи такие, как анализ больших данных, биоинформатика, требуют огромного количества ресурсов, что способствует росту числа гетерогенных распределенных систем
    \item Стандартные планировщики задач оптимизированы под гомогенные системы и не учитывают различия в производительности узлов и их доступность, что в гетерогенных системах приводит к простаиванию некоторых узлов или, наоборот, излишней загруженности отдельных
    \item Необходимость развития динамических алгоритмов балансировки в противовес сущестующим решениям с их статическими правилами
    \item Оптимизация всех вычислений, выполняющихся на платформе BOINC
\end{enumerate}

\subsection{Цели и задачи}
Целью данного проекта является разработка системы выдачи заданий для распределенной системы на платформе BOINC, в которой присутствуют сегменты кластеров. Для ее достижения необходимо было выполнить следующие задачи:
\begin{enumerate}
    \item Обзор систем балансировки нагрузки для кластеров и распределенных систем
    \item Выбор системы моделирования и адаптация модели под специфику гибридной распределенной системы
    \item Реализация/модификация системы балансировки нагрузки гибридной распределенной системы
    \item Тестирование системы балансировки на модели гибридной распределенной системы
\end{enumerate}

\subsection{Метрики оценки производительности}

Выбор метрик обусловлен необходимостью учета как общих показателей производительности, так и особенностей гетерогенной вычислительной среды.

\begin{itemize}
    \item Пропускная способность (Throughput) -- число выполненных задач за единицу времени. Формула: $T = \dfrac{N_{\text{completed}}}{\Delta t}$.
    \item Время выполнения задачи -- среднее, медиана, перцентили (95\%, 99\%). Для каждой задачи $i$: $R_i = t^{(i)}_{\text{finish}} - t^{(i)}_{\text{arrival}}$. Среднее: $\overline{R} = \dfrac{1}{N}\sum_i R_i$.
    \item Загруженность узлов -- для узла $j$: $u_j = \dfrac{\text{busy\_cpu\_time}_j}{\text{observation\_time}}$. Сводные метрики: среднее по узлам $\overline{u}$ и коэффициент вариации $CoV = \dfrac{\sigma(u_j)}{\overline{u}}$ (показывает балансировку).
    \item Общее время завершения набора задач -- $C = \max_{i}\{t_{\text{finish}}\} - \min_{i}\{t_{\text{arrival}}\}$
    \item Накладные расходы планировщика -- среднее время, затраченное на принятие решения по задаче, и доля от общего времени:
    \[
      OH_{\text{avg}} = \dfrac{1}{N}\sum_i t^{(i)}_{\text{sched}},\qquad
      OH_{\%} = \dfrac{\sum_i t^{(i)}_{\text{sched}}}{T_{\text{wall}}}\cdot 100\%.
    \]
\end{itemize}

Такой набор метрик позволяет оценить как производительность системы, так и качество балансировки нагрузки в условиях гетерогенной среды.

\subsection{Предполагаемые результаты проекта}
\begin{enumerate}
    \item Подробное описание методики моделирования гибридной распределенной системы
    \item Подробное описание системы балансировки нагрузки гибридной распределенной системы
    \item Анализ результатов моделирования системы
    \item Доклад о проекте на научной конференции
    \item Написание и публикация по итогам проделанной работы
\end{enumerate}

\section{Литературный обзор}
\subsection{Обзор предметной области}
В своем начале Интернет имел очень мало потребителей, так что любой сервер мог выдерживать нагрузку даже самого популярного сайта, к тому же его падение не стало бы большой проблемой. Со временем люди были вынуждены заняться вопросами масштабирования ввиду растущей нагрузки. И если сначала проблема решалась увеличением числа процессоров, их улучшением, увеличением объема оперативной памяти, то после стало понятно, что есть потолок такого вертикального масштабирования, так что появилось понятие балансировки нагрузки.~\cite{srv-lb}

В общем случае, порядок работы балансировщика такой:
\begin{enumerate}
    \item Получение входящих запросов от различных клиентов
    \item Расчет размера входящих запросов и формирование очереди запросов
    \item Периодическая проверка состояния серверов в пуле с помощью системы мониторинга
    \item Использование стратегии / алгоритма / эвристики для выбора подходящего сервера под каждый входящий запрос
\end{enumerate}

Сама идея балансировки нагрузки вносит в систему большое количество преимуществ:

\begin{enumerate}
    \item Упрощенное отслеживания трафика
    \item Улучшенное использование ресурсов
    \item Балансировка нагрузки на сеть на основании загруженности серверов
    \item Улучшенная доступность ресурсов
    \item Уменьшение избыточного потребления ресурсов
\end{enumerate}~\cite{lb-in-cloud}

Решения проблемы балансировки нагрузки делятся на 2 основных подхода в зависимости от того, принимает ли алгоритм решения на основе текущего состояния системы или нет: статический и динамический.

Статический подход предполагает наличие знаний о статусе системы, требованиях к ресурсам, времени связи и другим показателям. Балансировка при таком подходе достигается путем создания связей между списком процессоров и списком задач таким образом, чтобы достигалась максимальная производительность. Хотя назначения могут быть как детерминированными, так и вероятностными, текущее состояние системы не учитывается ни в одном из вариантов организации балансировки. Это существенный недостаток статических подходов, который оказывает существенное влияние на общую производительность из-за непредсказуемости колебаний нагрузки на систему. Некоторые методы, используемые при статической балансировке: перечисление и поиск в пространстве решений, теория графов, теория очередей.

При динамическом подходе, решения о балансировке нагрузки принимаются на основе текущего состояния системы; задачам разрешается динамически перемещаться с перегруженного узла на недостаточно загруженный узел для более быстрой обработки. Способность реагировать на изменения в системе является основным преимуществом динамического подхода к балансировке нагрузки. Хотя найти динамическое решение намного сложнее, чем статическое, динамическая балансировка нагрузки может повысить производительность за счет своей адаптивности. Далее сосредоточимся именно на таких подходах.

Механизм управления динамической балансировкой может быть распределенным, когда ответственность за балансировку распределяется между узлами, и нераспределенным, когда ответственность за балансировку лежит на одном или нескольких узлах; включает централизованный (один центральный узел управляет балансировкой всей системы) и полу распределенный (узлы делятся на кластеры с центральными узлами для балансировки внутри кластера) подходы.~\cite{lb-in-ds}

\subsection{Обзор существующих методов}
Рассмотрим методы балансировки нагрузки, применимые к распределённым вычислительным системам с акцентом на гибридные конфигурации (платформы типа BOINC + стабильные кластерные сегменты). Среди критериев для них можно выделить: простота мониторинга, устойчивость к churn (неконтролируемое подключение / отключение узлов), гетерогенность ресурсов, величина накладных расходов, масштабируемость. Всего можно выделить несколько групп подходов к балансировке нагрузки:

\begin{itemize}
    \item Классические статические методы (Round Robin, Random Strategy) — просты в мониторинге, однако не учитывают размер и время выполнения задач
    \item Классические динамические методы (Least-Connection) — также просты в мониторинге, но не учитывает общую нагрузку на узлы
    \item Адаптивные методы — балансировка происходит на основе метрик, вычисляемых в реальном времени; требуют постоянных вычислений
    \item Предиктивные ML подходы — планировщики, которые хорошо работают даже в сложных условиях, но требуют времени и данных для обучения модели~\cite{lb-adv}
    \item Существующие решения для BOINC (CluBoRun) — механизмы верификации, учета доступности, отслеживания churn и другие~\cite{boinc-lb}
\end{itemize}

Большинство классических и оптимизационных алгоритмов работают при допущении, что узлы доступны и их характеристики относительно стабильны. В средах типа BOINC это допущение неприменимо: изменчивость системы и разнородность устройств приводит к частому перераспределению задач и дополнительным затратам. Если в гомогенных системах перенос задач --- относительно стандартный процесс, в гибридных системах такой подход может привести к серьёзному росту накладных расходов из-за неоднородности устройств и отсутствия preemptive migration. Также нужно отметить, что для гибридных систем важно сохранять баланс между количеством собираемых метрик и их полезности для планировщика, что труднодостижимо при подходах с глобальным мониторингом узлов.~\cite{lb-in-cloud}

\subsection{Эвристики планирования задач}
Отдельно от сетевых балансировщиков в литературе по планированию задач (job scheduling) рассматриваются эвристики, непосредственно использованные в данной работе:
\begin{itemize}
    \item \textbf{LPT (Longest Processing Time first)} --- длинные задачи назначаются первыми; при параллельных однотипных машинах снижает makespan, но может ухудшить fairness и среднее время отклика коротких задач.
    \item \textbf{SJF (Shortest Job First)} --- короткие задачи получают приоритет; минимизирует среднее время отклика, но может голодать длинные workunits.
    \item \textbf{Random} --- нижняя граница качества планирования; полезен как контрольная политика.
\end{itemize}
Кластерные планировщики (Slurm, Torque) оптимизированы под гомогенные очереди с централизованным контролем и preemptive scheduling; они не учитывают pull-модель BOINC и добровольный churn узлов, поэтому не применимы «из коробки» к гибридной грид-системе.

\subsection{Особенности BOINC}
Платформа BOINC использует модель \emph{client-initiated}: вычислительный узел сам запрашивает работу через RPC. Scheduler сканирует shared-memory массив feeder, оценивает совместимость workunit с host и ранжирует кандидатов по score. Назначенная задача не мигрирует на другой host до timeout или повторной репликации. Поле \texttt{rsc\_fpops\_est} задаёт оценку сложности, которая может расходиться с реальным временем выполнения --- это важный источник ошибок для любой adaptive-политики. Решение CluBoRun~\cite{boinc-lb} интегрирует кластерные ресурсы в BOINC-проект, но не решает глобальную балансировку по классам устройств (кластер / desktop / mobile).

\subsection{Сравнение подходов}
Поймем, насколько применимы указанные подходы к гибридным системам и выделим их недочеты:
\begin{table}[h]
\centering
% \caption{Сравнение групп подходов к балансировке нагрузки в гибридных системах}
\footnotesize
\begin{tabular}{p{3.4cm} p{2.0cm} p{2.6cm} p{2.4cm} p{2.4cm}}
\hline
Группа методов & Тип & Требования мониторинга & Цель подхода & Минусы \\
\hline
Классические статические & Простые & Низкие & Простое планирование & Не учитывают размер задач \\
Классические динамические & Реактивные & Низкие/средние & Сократить число одновременных задач на узле & Не учитывает CPU-time узла \\
\hline
Адаптивные методы & Feedback-driven & Средние/высокие & Низкая задержка и оптимальная нагрузка на узлы & Высокие накладные расходы \\
\hline
Предиктивные ML-подходы & Обучаемые & Высокие & Минимизация целевой функции ошибки & Требуют данных и обучения \\
\hline
CluBORun & Интеграция кластера / устройств & Низкие/средние & Использование idle ресурсов, учёт доступности & Не решает глобальную балансировку \\
\hline
\end{tabular}
\end{table}

\subsection{Выводы из обзора}
\begin{itemize}
    \item На данный момент существует широкий набор методов балансировки, однако немногие работы систематически исследуют их поведение в гибридных системах
    \item Основные пробелы: недостаточная оценка влияния churn, упрощённое моделирование сетевых затрат при перераспределении задач, отсутствие оценки накладных расходов мониторинга и планировщика
\end{itemize}

Данная работа направлена на устранение выявленных ограничений за счёт оптимизации накладных расходов планировщика, более точного учёта динамической доступности узлов, снижения частоты миграций задач и повышения общей производительности системы при рациональном использовании вычислительных ресурсов.

\section{Модель гибридной распределённой системы}

\subsection{Выбор системы моделирования}
Для экспериментальной проверки политик балансировки выбран локальный стенд на базе Docker Compose и реального scheduler BOINC~8.3.0. Альтернативы --- discrete-event симулятор или GPU-стенд --- отвергнуты: симулятор требует отдельной калибровки под протокол BOINC; GPU-модель на Apple Silicon не даёт корректного passthrough в Linux-контейнеры и смещает фокус с CPU-centric volunteer computing.

Основные преимущества выбранного подхода:
\begin{itemize}
    \item полный RPC-цикл BOINC (feeder, feasibility checks, replication);
    \item воспроизводимость через скрипты и \texttt{manifest.json};
    \item контролируемая гетерогенность через Docker cgroup и BOINC preferences;
    \item ground truth по времени выполнения через приложение \texttt{uppercase} с параметром \texttt{--cpu\_time}.
\end{itemize}

\subsection{Топология стенда}
Стенд моделирует гибридную систему из шести вычислительных hosts:
\begin{itemize}
    \item \texttt{client-cluster}, \texttt{client-cluster-2} --- сегменты кластера (4 CPU, 4~GiB);
    \item \texttt{client-desktop} --- персональный компьютер (2 CPU, 2~GiB);
    \item \texttt{client-low-power} --- маломощный узел (1 CPU, 1~GiB);
    \item \texttt{client-phone}, \texttt{client-phone-2} --- мобильные профили (0.5 CPU, 512~MiB).
\end{itemize}
Каждый контейнер имеет отдельный BOINC data directory и регистрируется как независимый host. Инфраструктура включает MySQL, Apache (scheduler CGI) и сервис сборки проекта.

\subsection{Workload и профили churn}
Основной набор --- 36 задач: 18 short (5~s CPU), 12 medium (30~s), 6 long (120~s). Для коротких smoke-прогонов использовалось 12 задач. Профили доступности:
\begin{itemize}
    \item \textbf{stable} --- все узлы online весь прогон;
    \item \textbf{heavy} --- у каждого клиента несколько циклов \texttt{pause/unpause} с рассинхронизацией (\texttt{CHURN\_STAGGER\_SECONDS}).
\end{itemize}
Churn реализован через паузу контейнера (узел пропадает из сети, но BOINC-процесс не перезапускается). После валидации всех workunits churn останавливается.

\subsection{Ограничения локального воспроизведения}
\begin{itemize}
    \item масштаб --- 6 hosts на одной машине, упрощённая сеть;
    \item BOINC часто сообщает одинаковый \texttt{p\_fpops} для всех hosts --- в коде добавлен \texttt{host\_speed\_factor} по hostname;
    \item cgroup CPU $\neq$ реальное железо (архитектура, энергосбережение);
    \item \texttt{pause/unpause} мягче реального churn (переустановка клиента, смена IP);
    \item нет preemptive migration задач --- как в production BOINC;
    \item \texttt{max\_wus\_to\_send=1}, \texttt{max\_wus\_in\_progress=2} --- искусственное ограничение захвата очереди быстрыми клиентами.
\end{itemize}
Эти ограничения не отменяют ценность A/B-сравнения политик на одном стенде, но не позволяют напрямую экстраполировать абсолютные значения метрик на production BOINC.

\section{Реализация системы балансировки}

\subsection{Точка интеграции}
Модификация выполнена в модуле \texttt{sched\_send.cpp}: при \texttt{custom\_load\_balancer=1} вызывается \texttt{send\_work\_custom()} вместо штатного \texttt{send\_work\_score()}. Feature flag в \texttt{config.xml} позволяет проводить A/B-эксперименты без пересборки бинарника.

\subsection{Архитектурный принцип}
Функция \texttt{send\_work\_custom()} сохраняет весь контракт штатного scheduler: сканирование feeder shared memory, feasibility checks (deadline, память, квоты, replication), locking и обновление БД. \textbf{Изменяется только ranking} допустимых кандидатов --- минимально инвазивная модификация, снижающая риск нарушить штатное поведение BOINC.

\subsection{Экспериментальные политики}
\begin{table}[h]
\centering
\caption{Сравнение политик dispatch}
\footnotesize
\begin{tabular}{p{2.2cm} p{1.6cm} p{8.5cm}}
\toprule
Политика & Flag & Суть ranking \\
\midrule
baseline & \texttt{custom\_load\_balancer=0} & Штатный BOINC score: reliability, locality, beta, estimated duration \\
random & 1, \texttt{random} & BOINC score + детерминированный hash(\texttt{result\_id}, \texttt{host\_id}) \\
lpt & 1, \texttt{lpt} & BOINC score + $\log(\text{predicted\_runtime})$ --- длинные задачи первыми \\
sjf & 1, \texttt{sjf} & BOINC score $-$ $\log(\text{predicted\_runtime})$ --- короткие первыми \\
hybrid & 1, \texttt{hybrid} & Мульти-критериальный score (см. ниже) \\
\bottomrule
\end{tabular}
\end{table}

Политики random, lpt и sjf --- \textbf{контрольные}: они изолируют отдельные гипотезы (случайное назначение, минимизация makespan, минимизация response time) и не претендуют на production-решение.

\subsection{Hybrid-политика}
Для пары (workunit $i$, host $j$) вычисляется:
\begin{align*}
\hat{t}_{ij} &= \frac{\text{estimate\_duration}(i,j)}{\text{host\_speed\_factor}(j)}, \\
\text{size\_affinity}_{ij} &= -\left|\log\frac{\hat{t}_{ij}}{T_{\text{target}}}\right|, \\
\text{deadline\_press}_i &= \frac{T_{\text{target}}}{\text{delay\_bound}_i}, \\
\text{runtime\_penalty}_{ij} &= \log\left(1 + \frac{\hat{t}_{ij}}{T_{\text{target}}}\right), \\
\text{host\_fit}_{ij} &= \log(1 + v_j) - \log(1 + \hat{t}_{ij}), \\
\text{score}_{ij} &= 25 \cdot s^{\text{BOINC}}_{ij}
  + w_{\text{size}} \cdot \text{size\_affinity}_{ij}
  + w_{\text{dl}} \cdot \text{deadline\_press}_i \\
  &\quad - w_{\text{rt}} \cdot \text{runtime\_penalty}_{ij}
  + 0.5 \cdot w_{\text{size}} \cdot \text{host\_fit}_{ij}.
\end{align*}
Здесь $s^{\text{BOINC}}_{ij}$ --- штатный BOINC score, $v_j$ --- коэффициент мощности host (cluster=4, desktop=2, low-power=1, phone=0.5), $T_{\text{target}}$ --- целевое время выполнения (30~s).

\textbf{Теоретическое обоснование:}
\begin{itemize}
    \item \texttt{size\_affinity} и \texttt{host\_fit} сопоставляют мощность host и длительность задачи;
    \item \texttt{deadline\_press} предотвращает голодание срочных workunits;
    \item \texttt{runtime\_penalty} снижает назначение длинных задач на медленные hosts;
    \item множитель $25\times$ (вместо $100\times$ у контрольных политик) сохраняет приоритет reliability/locality, но даёт custom-членам влиять на порядок кандидатов.
\end{itemize}

\section{Методика экспериментов}

\subsection{Протокол}
Эксперименты запускались скриптом \texttt{experiments/run\_experiment.sh all}: подряд выполняются пять политик с одним \texttt{PAIR\_ID}. Для каждой политики создаётся свежая БД при одинаковом workload и профиле churn. По завершении экспортируются SQL-данные, рассчитываются метрики и сохраняются в \texttt{results/<pair-id>\_<policy>/}.

\subsection{Группы прогонов}
Проведено девять прогонов режима \texttt{all} (по три повтора в каждой группе):
\begin{table}[h]
\centering
\caption{Группы экспериментальных прогонов}
\footnotesize
\begin{tabular}{lccc}
\toprule
Группа & \texttt{pair\_id} & Churn & Jobs \\
\midrule
A (основная) & 20260914T173104Z, 191425Z, 194153Z & stable & 36 \\
B (exploratory) & 20260914T213013Z, 073506Z, 083656Z & heavy & 12 \\
C (stress) & 20260915T132706Z, 140620Z, 160921Z & heavy & 36 \\
\bottomrule
\end{tabular}
\end{table}

\subsection{Сравниваемые политики}
Внутри каждого прогона сравниваются baseline (штатный BOINC), random, lpt, sjf и hybrid при идентичных условиях. Значения в таблицах ниже --- \textbf{среднее арифметическое} по трём повторам группы. Во всех прогонах \texttt{deadline\_miss\_rate} = 0.

\section{Результаты и анализ}

\subsection{Stable, 36 задач --- основной сценарий}
\begin{table}[h]
\centering
\caption{Средние метрики: stable churn, 36 задач}
\footnotesize
\begin{tabular}{lrrrrrr}
\toprule
Политика & $T$ & $\overline{R}$, s & $R_{95}$, s & Makespan, s & Load CV & Jain \\
\midrule
baseline & 0.121 & 130.3 & 183.9 & 251.7 & 0.501 & 0.773 \\
random   & 0.125 & 120.2 & 197.9 & 248.7 & 0.463 & 0.750 \\
lpt      & \textbf{0.126} & 139.5 & 201.3 & \textbf{238.3} & \textbf{0.334} & 0.615 \\
sjf      & 0.117 & \textbf{116.0} & 209.7 & 261.0 & 0.588 & \textbf{0.778} \\
hybrid   & 0.117 & 136.0 & 214.8 & 246.0 & 0.449 & 0.736 \\
\bottomrule
\end{tabular}
\end{table}

\textbf{LPT} показывает лучшие makespan ($\approx$238~s), throughput и баланс нагрузки (Load CV $\approx$ 0.33): длинные задачи сразу уходят на быстрые cluster-hosts, сокращая хвост набора. Цена --- худший Jain ($\approx$0.62): cluster-узлы получают непропорционально больше workunits.

\textbf{Random} даёт лучший mean response ($\approx$120~s) и высокий throughput: при малом числе hosts отсутствие систематического bias иногда выигрывает у «умного» неправильного предсказания.

\textbf{SJF} минимизирует response time и максимизирует fairness, но ухудшает makespan ($\approx$261~s) --- long workunits ждут в хвосте.

\textbf{Baseline} --- сильный all-rounder без явных провалов; штатный score уже учитывает reliability и estimated duration.

\textbf{Hybrid} на stable-36 \textbf{не превосходит baseline} по throughput и makespan; улучшает Load CV относительно baseline, но уступает LPT.

\subsection{Heavy churn, 12 задач}
\begin{table}[h]
\centering
\caption{Средние метрики: heavy churn, 12 задач}
\footnotesize
\begin{tabular}{lrrrrr}
\toprule
Политика & $T$ & $\overline{R}$, s & Makespan, s & Load CV & Jain \\
\midrule
baseline & 0.046 & 91.1 & 201.0 & 0.943 & 0.816 \\
random   & \textbf{0.049} & 96.7 & \textbf{189.0} & 0.810 & 0.699 \\
lpt      & 0.047 & 107.1 & 204.3 & 0.758 & 0.816 \\
sjf      & 0.049 & \textbf{84.9} & 190.0 & 0.908 & 0.752 \\
hybrid   & 0.045 & 104.0 & 205.7 & \textbf{0.676} & \textbf{0.838} \\
\bottomrule
\end{tabular}
\end{table}

При малом числе задач наблюдается высокая дисперсия; абсолютные значения несопоставимы с 36-job группами. \textbf{Hybrid} лидирует по Jain и Load CV, но проигрывает по throughput. Группа B годится как иллюстрация влияния churn, но не как основное доказательство.

\subsection{Heavy churn, 36 задач}
\begin{table}[h]
\centering
\caption{Средние метрики: heavy churn, 36 задач}
\footnotesize
\begin{tabular}{lrrrrrr}
\toprule
Политика & $T$ & $\overline{R}$, s & $R_{95}$, s & Makespan, s & Load CV & Jain \\
\midrule
baseline & \textbf{0.097} & 157.1 & 233.5 & 324.3 & 0.558 & 0.787 \\
random   & 0.095 & 152.5 & 235.1 & 324.0 & 0.471 & 0.782 \\
lpt      & 0.095 & 166.2 & \textbf{215.8} & \textbf{261.3} & \textbf{0.263} & 0.620 \\
sjf      & 0.101 & \textbf{134.7} & 235.0 & 302.0 & 0.379 & 0.817 \\
hybrid   & 0.088 & 154.5 & 227.3 & 349.7 & 0.599 & \textbf{0.923} \\
\bottomrule
\end{tabular}
\end{table}

Сценарий наиболее стрессовый и нестабильный (высокий разброс между прогонами). \textbf{LPT} минимизирует makespan и Load CV, когда «работает»; в отдельных прогонах throughput падал из-за совмещения long jobs с offline cluster-hosts.

\textbf{Hybrid} стабильно даёт лучший Jain ($\approx$0.923) --- почти равномерное распределение завершённых задач по шести hosts; это следствие \texttt{host\_fit} и \texttt{size\_affinity}. Однако hybrid проигрывает по makespan ($\approx$350~s vs $\approx$324~s у baseline) и throughput: при heavy churn медленные hosts получают workunits, пока быстрые offline.

\subsection{Сводный trade-off}
\begin{table}[h]
\centering
\caption{Качественное сравнение политик по сценариям}
\footnotesize
\begin{tabular}{lcc}
\toprule
Политика & stable-36 & heavy-36 \\
\midrule
LPT      & скорость $\checkmark$, fairness $\times$ & скорость $\sim$, fairness $\times$ \\
Hybrid   & скорость $\sim$, fairness $\sim$       & скорость $\times$, fairness $\checkmark\checkmark$ \\
Baseline & скорость $\checkmark$, fairness $\checkmark$ & скорость $\checkmark$, fairness $\checkmark$ \\
\bottomrule
\end{tabular}
\end{table}

\subsection{Выводы по экспериментам}
\begin{enumerate}
    \item Разработан и интегрирован кастомный dispatch layer для BOINC с пятью режимами сравнения.
    \item На stable-36 \textbf{LPT} лучше по makespan, throughput и балансу нагрузки; \textbf{baseline} --- надёжная база без провалов.
    \item \textbf{Hybrid} улучшает fairness при churn, но ухудшает throughput/makespan из-за назначения работы ослабленным hosts в моменты их кратковременной доступности.
    \item Предложенная политика --- проверяемая гипотеза, а не гарантированное улучшение; требуется настройка весов или учёт availability.
\end{enumerate}

\section{Заключение}

В рамках курсовой работы выполнены все поставленные задачи:
\begin{enumerate}
    \item проведён обзор методов балансировки для кластеров и распределённых систем с акцентом на BOINC и гибридные конфигурации;
    \item выбрана и адаптирована CPU-модель гибридной системы на Docker/BOINC с шестью профилями hosts и управляемым churn;
    \item реализована модификация \texttt{sched\_send} с политиками baseline, random, lpt, sjf и hybrid;
    \item проведено девять сравнительных прогонов; основной вывод --- LPT минимизирует makespan на stable-стенде, hybrid улучшает fairness при churn ценой скорости.
\end{enumerate}

\textbf{Ограничения:} малый масштаб (6 hosts), упрощённый churn, hostname-based speed hints, отсутствие GPU plan class.

\textbf{Направления дальнейшей работы:} учёт битовой availability в score, segment-aware routing, sensitivity analysis весов hybrid-политики, per-class breakdown (short/medium/long) по типам hosts.







% Существует множество типовых методов балансировки нагрузки между узлами. Классические методы, как несложно понять, в большинстве своем статические, то есть не учитывают текущее состояние системы, однако есть и несколько динамических, пытающихся оптимизировать распределение задач между узлами.

% Least-Connection Load Balancing -- очередной запрос идет на узел, который в данный момент обрабатывает наименьшее число запросов

% Round-Robin Load Balancing -- если в системе n узлов, то запросы будут распределяться следующим образом: узел 1, узел 2, …, узел n, узел 1 и так далее

% Random Strategy based Load Balancing -- узел выбирается случайно

% Clustering based Load Balancing -- данный подход заключается в преобразовании системы разнородных узлов в группы однородных, что позволяет сбалансировать узлы с помощью ранее рассмотренных методов. Алгоритм кластеризации выполняется не централизованно, а каждым узлом, который может принимать решения о “создании” или “отключении” связи от соседнего узла. Балансировка нагрузки в данном случае есть итеративное выполнение шагов: случайный узел выбирает себя инициатором и выбирает узел-посредник; посредник выбирает соседний узел, который сопоставим с инициатором, создает связь между этими двумя узлами, а затем отключается. Таким образом, группа однородных узлов объединяется в гетерогенную систему, которая затем может распределять нагрузку внутри себя.~\cite{lb-in-cloud}

% В последнее время наиболее популярными становятся алгоритмы балансировки, применяющие машинное обучение. Вот несколько алгоритмов, предложенные в последние годы:

% Adaptive Load Balancer -- предполагает непрерывный мониторинг показателей – время отклика узла, текущая нагрузка, задержка сети, использование ресурсов. Он динамически подстраивается под условия реального времени, перенаправляя нагрузку с загруженных или долго отвечающих узлов на менее загруженные с целью минимизации задержки и предотвращения перегрузки серверов, обеспечивая эффективное использование ресурсов и повышая производительность.

% Predictive Load Balancer -- основан на анализе исторических данных о трафике, производительности серверов и состоянии сети, выявляя тенденции и закономерности, которые могут быть использованы в будущем при балансировке нагрузки. Используя методы машинного обучения, алгоритм генерирует прогнозы будущего трафика и нагрузки на сервер, которые основаны на выявленных закономерностях и показателях, получаемых в режиме реального времени, и корректирует на основе этого распределение запросов. Алгоритм призван прогнозировать и устранять потенциальные проблемы с производительностью до их возникновения.~\cite{lb-adv}

% \subsection{Сравнение существующих подходов}
% \begin{table}[h] \centering \caption{Сравнение классов методов балансировки} \begin{tabular}{p{3.5cm} p{2cm} p{2.5cm} p{2.5cm} p{2.5cm}} \hline Метод & Тип & Требования к мониторингу & Главная цель & Ключевые ограничения \ \hline Round-Robin / простые эвристики & Централиз./Детермин. & Низкие & Простота и низкий overhead & Игнорируют размер/нагрузку задач \ Global heuristics (bin-packing, min-makespan) & Централиз. & Высокие (глобальная информация) & Минимизация makespan & Высокий вычислительный и мониторинговый overhead \ Adaptive load-aware & Гибрид/динамич. & Средние & Балансировка latency и utilisation & Зависимость от свежести данных \ ML-based & Централиз./обучаемые & Высокие (истор. данные) & Оптимизация сложных целей & Требуется обучение/устойчивость к дрейфу \ \hline \end{tabular} \end{table}
% \begin{table}[]
%     \centering
%     \begin{tabular}{c|c}
%          &  \\
%          & 
%     \end{tabular}
%     \caption{Существующие методы балансировки}
%     \label{tab:placeholder}
% \end{table}

	
\newpage 
\printbibliography[heading=bibintoc] 

% \begin{thebibliography}{0}
% 	\bibitem{chirkova18}\hypertarget{chirkova18}{}
% 	\href{https://arxiv.org/abs/1810.10927}
% 	{Nadezhda Chirkova, Ekaterina Lobacheva, Dmitry Vetrov. Bayesian Compression for Natural Language Processing. In EMNLP 2018.}
% \end{thebibliography}
	
\newpage
\appendix

% \section{Пример секции аппендикса}

\end{document}
